% !TeX program = pdflatex
% papers/medida.tex — a teoria da medida deste sistema, derivada do corpo estelar.
% Auto-contido: não depende de teoria.tex nem de catalogo.tex para ser lido.
% Compila sozinho:  pdflatex medida.tex
\documentclass[11pt,a4paper]{article}
\usepackage[utf8]{inputenc}\usepackage[T1]{fontenc}\usepackage[portuguese]{babel}
\usepackage{amsmath,amssymb,amsthm}
\usepackage[margin=2.3cm]{geometry}
\usepackage{booktabs}\usepackage{xcolor}
\usepackage[hidelinks]{hyperref}

\definecolor{cinza}{gray}{0.35}
\newcommand{\code}[1]{\texttt{\small #1}}
\newcommand{\medido}[1]{\par\medskip\noindent{\small\color{cinza}\textbf{Medido.} #1}\par\medskip}
\newtheorem{teorema}{Teorema}
\newtheorem{definicao}[teorema]{Definição}
\newtheorem{corolario}[teorema]{Corolário}
\renewcommand{\proofname}{Demonstração}
\newcommand{\dg}{^{\dagger}}

\title{\textbf{A medida}\\[3pt]
  {\large saltar sobre o irracional, furar o que não conta}}
\author{Aarão Melo Lopes}
\date{8 de agosto de 2026}

\begin{document}
\maketitle

\begin{abstract}
\noindent Este documento funda a medida deste sistema, derivada do \emph{corpo estelar} e da mesma
exigência única. A tese cabe em três verbos, e nenhum deles é \emph{aproximar}:
\textbf{contar} --- a medida é o contador do relógio lido em por-unidade, e $\pi$ já sai por
contagem de pontos inteiros; \textbf{saltar} --- um irracional não se atinge e não é preciso:
encaixa-se pelo \emph{par} que o encaixota, uma coordenada por baixo e uma por cima, e a forma
decimal é \emph{meia} medida com o nome do par; \textbf{furar} --- um ponto custa menos que
qualquer régua, uma família contável fura-se com custo total tão pequeno quanto se queira, e por
isso os racionais \emph{não pesam}: fura-se e a régua não se move.

\smallskip
\noindent Daí deriva-se a integral que o clássico atribui a Lebesgue, e a derivação é uma
involução já provada: \textbf{cortar a imagem é o dual de cortar o domínio} --- $f\leftrightarrow
f^{-1}$, a Lei~2 no nível um --- e a conservação $\int f+\int f^{-1}=b\,f(b)-a\,f(a)$, já medida
nos metais, é exactamente a troca de um corte pelo outro. A função que vale $1$ nos racionais
integra $0$ numa linha: furam-se os racionais e sobra a régua inteira.

\smallskip
\noindent E a medida tem \textbf{dois objectos}, porque a cruz tem duas coordenadas: o
\emph{conjunto}, cujo tamanho a soma mede --- contar, saltar, furar ---, e o \emph{corpo}, cuja
\textbf{espécie} o produto ordena --- e «espécie» é a \emph{régua} com mais um nome (assinatura,
cifra: o segundo grupo de nomes da \emph{Teoria}), lida junto com a \textbf{dinâmica temporal},
porque o espectro tem duas colunas e o crânio tem dois ossos. Medir a espécie célula a célula \emph{não mede} --- não por
imprecisão: por categoria. A espécie classifica-se pelo crânio, e crânios são meia dúzia: a
assinatura, inteiros contados pelo relógio em por-unidade. É por isso que Fourier não mede os
corpos daqui --- avalia no círculo e as raízes do metal não estão no círculo --- e a transformada
que mede é a \textbf{universal}: avaliação nas raízes da própria borda, com a \emph{inversa} a
vestir a roupa da espiral no fim, sem sobra e sem excesso.

\smallskip
\noindent E fecha-se com o que isto explica do próprio projecto: \emph{todos} os objectos desta
teoria --- os metais, os corpos $K_{n,m}$, os Pisot --- têm medida nula, e é por isso que a análise
de \emph{quase todo} número não vê nenhum deles. A medida não é um capítulo a mais: \textbf{é o
mapa da cegueira do clássico}, e este projecto vive inteiro dentro dela.
\end{abstract}

\section{A exigência, herdada}

O \emph{corpo estelar} fixa como se mede aqui, e esta teoria não acrescenta uma segunda doutrina
--- realiza a primeira:

\begin{quote}
\textbf{A medida não compara: reverte, e lê o resíduo.} E mede-se \emph{pela metade}: o resíduo
onde \emph{tem} de ser zero diz que \textbf{existe}; o resíduo onde \emph{não pode} ser zero diz
que é \textbf{único}.
\end{quote}

\noindent O que este documento faz é aplicar essa doutrina ao objecto que faltava: \emph{o
tamanho de um conjunto}. E a resposta vai contra o hábito de um século de sala de aula: tamanho
não se aproxima por baixo, dígito a dígito --- \textbf{conta-se por cima, salta-se o que não se
atinge, e fura-se o que não conta}.

\section{A mineração: o que já estava derivado}

Antes de escrever, procurou-se --- a regra do projecto. Três peças já existiam, e este documento
cita-as em vez de as refazer:

\begin{center}\small
\begin{tabular}{@{}lll@{}}
\toprule
a peça & onde vive & o que dá a esta teoria \\
\midrule
o \textbf{contador} $\lfloor t/d\rfloor+1$ & a \emph{Teoria}, o relógio & a medida é contagem por escala \\
o \textbf{pente} autodual & a \emph{Teoria}, a Dirac & as marcas, e o dual de passo $N/d$ \\
os \textbf{encaixantes} & a \emph{Teoria} & o salto: o par que encaixota o ponto \\
$\pi$ \textbf{por contagem} & \code{tests/pimonte.c} & contar bate tentar, medido \\
a \textbf{medida nula} dos objectos & a \emph{Teoria}, conclusão & o mapa da cegueira do clássico \\
o corpo \textbf{métrico} & o \emph{Catálogo} (Fisher) & \emph{não} é isto: mede distância, não tamanho \\
\bottomrule
\end{tabular}
\end{center}

\noindent A verificação pedida fecha assim: o \emph{corpo métrico} do catálogo é a régua de
Fisher --- $g(p)=1/\bigl(p(1-p)\bigr)$, $(1-s^2)\,g(p)=4$ --- e mede \emph{distância entre
distribuições}, não tamanho de conjuntos; o \emph{relógio} da teoria tem o contador e o pente, que
são a matéria-prima daqui, mas não tem a medida de um conjunto qualquer nem a integral. \textbf{A
medida não estava derivada; as suas peças estavam.} Este documento é a montagem.

\medido{o contador foi contado \emph{marca a marca} contra a forma fechada em $18\,030$ pares
(passo, instante), dois caminhos independentes sem uma discordância, com o enquadramento exacto
$t<\text{contador}\cdot d\leq t+d$ em toda a amostra (\code{tests/relogio.c}); o dual do pente de
passo $d$ é o pente de passo $N/d$ em $8899$ casos de $8899$, por aritmética modular pura; e a fase
$p/q$ dá exactamente $q$ marcas em $12\,232$ fases primitivas, tudo em $\mathbb{Z}/q$, sem um
arredondamento a decidir coincidência (\code{tests/fase\_regua.c}).}

\section{A definição: a medida é o contador, em por-unidade}\label{sec:def}

\begin{definicao}[a medida de um intervalo]
Num relógio de passo $1/q$, o \textbf{peso} de um intervalo $I$ é o número de marcas que caem
nele --- o contador. A \textbf{medida} é o peso em por-unidade,
\[
  \lambda(I) \;=\; \frac{\#\{\text{marcas de passo }1/q\text{ em }I\}}{q},
\]
lida na escala em que o passo deixa de distinguir --- e o enquadramento do contador diz que duas
escalas quaisquer concordam a menos de \emph{um} passo, logo a leitura estabiliza e é única.
\end{definicao}

\noindent \emph{Nenhum limite de somas de Riemann entrou}: entrou o contador, que já estava
derivado, e a divisão por $q$, que é o por-unidade de sempre. Para $I=[a,b]$ o enquadramento dá
$\lambda(I)=b-a$ --- \textbf{o comprimento é o caso mais simples da contagem}, e não o contrário.

\smallskip
\noindent E a medida é \emph{leitura}, não movimento: é a coordenada da cruz que \textbf{mede} ---
aditiva, cega ao sinal, a Lei~1. A coordenada que \textbf{ordena} --- multiplicativa, a Lei~2 ---
é a outra metade, e vai aparecer no~\S\ref{sec:lebesgue} como a involução que troca os cortes.
\emph{Uma sem a outra é meia medida}, e é exactamente o defeito da forma decimal, já a seguir.

\medido{$\pi$ sai da definição acima com $N(R)/R^{2}$ --- pontos do reticulado no disco, tudo em
inteiros: com $R=3000$ o desvio é de $15$ milionésimos, contra $28\,408$ em $R=10$ --- apertar a
escala mede melhor e a convergência \emph{oscila} em vez de descer sempre. E o controlo é o corpo
térmico: o Monte Carlo --- proposta cega mais critério dentro/fora --- chega ao mesmo número e
chega \emph{pior}, $84$ contra $15$. \textbf{Quem tem a régua ordenada não precisa de tentar}
(\code{tests/pimonte.c}).}

\section{Saltar: o irracional não se aproxima --- encaixa-se}\label{sec:saltar}

Dois factos já provados na \emph{Teoria} decidem isto, e nenhum é deste documento:

\begin{enumerate}
\item \textbf{nenhum passo próprio atinge o alvo} --- a órbita racional aproxima-se de $\pi$ sem
chegar, e todo ponto fixo de operador unimodular é quadrático: \emph{o transcendente fica de fora
por consequência, não por escolha};
\item \textbf{o encaixante tem ponto, e é único} --- uma sucessão de encaixados não vazios tem
intersecção, e com a medida a ir a zero a intersecção é \emph{um} ponto.
\end{enumerate}

\noindent Juntos dizem o que fazer com um irracional: \emph{não ir lá}. O encaixante é a cruz de
uma órbita que não fecha --- \textbf{uma coordenada por baixo e uma por cima}, o par que o
encaixota --- e a medida do encaixante lê-se pelo contador em cada escala, \emph{sem nunca tocar o
ponto}. Saltar sobre ele é isto: o par passa por cima, e o que fica medido é o intervalo, não o
dígito.

\smallskip
\noindent \textbf{E a forma decimal é meia medida com o nome do par.} Escrever $\pi=3{,}14159\ldots$
é dar só a coordenada de baixo --- cada truncamento aproxima por \emph{um} lado, e quem só tem um
lado não sabe a que distância está. O par de convergentes da fracção contínua tem as duas: encosta
por baixo \emph{e} por cima, alternando, e o erro fica \emph{enquadrado} em vez de estimado. É o
mesmo defeito que o corpo estelar aponta na asserção que só olha um resíduo: \emph{estar pela
metade é o que dissipa}. Aproximar o infinito dígito a dígito é pagar infinitos passos pela metade
da informação; \textbf{encaixar paga dois números por escala e leva a informação inteira}.

\medido{a órbita racional aproxima-se sem chegar: dos $2\,376$ pontos racionais do círculo, os
$400$ que realizam o alvo são todos degenerados, e o traço mais próximo alcançado é
$-3968/1985$, estritamente acima do alvo $-2$; e $\pi$ é produzido por três somas alternadas
independentes em aritmética \emph{inteira} que concordam nas $24$ primeiras casas --- nenhum
dígito escrito à mão (\code{tests/pidual.c}). A prova dos nove fecha em $23\,200$ pares com o erro
de uma unidade a acusar em todos (\code{tests/reversao.c}).}

\section{Furar: o que não conta, não pesa}\label{sec:furar}

\begin{teorema}[um ponto tem medida nula]
Em toda escala $q$, um ponto cabe num encaixante com \emph{uma} marca, de peso $1/q$. Como a
leitura vale em toda escala e $1/q$ desce sem parar, a medida do ponto é menor que toda régua ---
e o único valor menor que toda régua é \textbf{zero}.
\end{teorema}

\begin{teorema}[fura-se uma família contável, e a soma não aparece]\label{thm:furos}
Sejam $x_1,x_2,\ldots$ pontos, um por natural. Cubra-se o $k$-ésimo com um encaixante de medida
$\varepsilon/2^{k}$. A soma dos custos é
\[
  \sum_{k\geq1}\frac{\varepsilon}{2^{k}} \;=\; \varepsilon ,
\]
e $\varepsilon$ é qualquer. \emph{Logo o conjunto inteiro tem medida nula.} E o $2^{k}$ não veio
de fora: é a torre --- $\dim A_{n+1}=2\dim A_n$ --- descida em vez de subida, \textbf{a mesma
duplicação lida como metade}.
\end{teorema}

\begin{proof}
A medida é contagem em por-unidade e o contador é monótono: o peso da união não excede a soma dos
pesos, escala a escala. A soma geométrica fecha em $\varepsilon$ dentro do anel --- é a série da
metade, e a metade é exacta aqui porque a divisão por dois é a da promoção, sempre inteira no par.
\end{proof}

\noindent \textbf{E é isto que o projecto pedia desde o início: fura alguns pontos e já era.} Os
racionais são contáveis --- a cabeça em $\mathbb{Z}$, a palavra em $\mathbb{N}^{*}$, a
correspondência já está no catálogo --- logo \textbf{$\mathbb{Q}$ fura-se inteiro e a régua não se
move}: $\lambda([a,b]\setminus\mathbb{Q})=b-a$. Não se aproximou irracional nenhum para o dizer;
furou-se o que conta de outra maneira.

\smallskip
\noindent \textbf{E daí o mapa da cegueira.} A conclusão da \emph{Teoria} já o dizia: os corpos
$K_{n,m}$, os metais, os Pisot --- \emph{todos} de medida nula, todos invisíveis a Gauss--Kuzmin,
Khinchin e Lévy, que legislam sobre \emph{quase todo} número e sobre nenhum em particular. Agora a
frase tem o mecanismo: \textbf{quase todo é o complementar de um furo}, e este projecto vive
inteiro \emph{dentro} dos furos. A medida não diminui os objectos daqui --- diz por que o clássico
não os vê, e di-lo com uma contagem.

\section{O bit aleatório enche o espaço: a medida por Monte Carlo--Metropolis}\label{sec:metropolis}

A \emph{Teoria} mostra que um único bit \emph{determinista} --- a unidade $b\in\{0,1\}$ sob a
convolução (o gato, $\times\sigma$) --- \textbf{desenha} uma estrutura: um rasto de dimensão um. Mas
uma curva, pela secção anterior, tem \textbf{medida nula}: um só bit determinista traça o esqueleto e
não \emph{enche} nada. Para a medida --- a área, o contínuo cheio --- falta o outro lado do bit: o
\textbf{acaso}.

\begin{teorema}[um bit aleatório enche o espaço, e o rasto é a medida]\label{thm:metropolis}
Um único bit sob um passeio de \textbf{Metropolis} --- proposta $\pm1$ numa direcção sorteada, aceite
pela simetria --- é \emph{ergódico}: o seu rasto visita \textbf{todas} as casas do reticulado, e a
fracção do tempo que ele passa num conjunto $S$ converge para a medida de Lebesgue $\lambda(S)$. E
\textbf{não interessa qual bit}: a ergodicidade apaga a semente e a posição inicial --- qualquer bit,
em qualquer \emph{thread}, enche o mesmo espaço.
\end{teorema}

\noindent \textbf{E é o dual da contagem, não a sua ruína.} A régua ordenada (\S\ref{sec:def}, o
contador) mede o disco \emph{exacto}, contando as casas dentro; o bit aleatório mede-o pela fracção do
rasto --- Monte Carlo. São \textbf{o par}: um é discreto e exacto, o outro é contínuo e enche, e os
dois dão a mesma $\lambda$. \emph{A convergência é medida, não assumida}: o erro médio sobre muitas
sementes desce com os passos (a lei dos grandes números), e é isso que faz do acaso uma régua e não um
palpite.

\begin{center}\small
\begin{tabular}{@{}lll@{}}
\toprule
o bit & o que faz & a medida \\
\midrule
\textbf{determinista} (o gato) & \emph{desenha} uma curva & \textbf{zero} (um rasto de dimensão um) \\
\textbf{aleatório} (Metropolis) & \emph{enche} a área & \textbf{cheia} (a fracção do rasto é $\lambda$) \\
\bottomrule
\end{tabular}
\end{center}

\noindent \textbf{O mesmo bit, os dois lados.} Desenhar ou encher --- a curva de medida zero e a área
de medida cheia --- são a estaca e a cruz do bit: um traça o que \emph{é}, o outro mede \emph{onde
está}. \emph{Aqui} o rasto determinista do gato é uma curva simples, e uma curva simples não pesa; o
acaso é que enche. \textbf{Mas encher não é privilégio do acaso}: há uma trajectória \emph{determinista}
que também enche --- a \textbf{curva de Peano}, 1D e de imagem 2D. A diferença entre as duas não é
encher: é a \emph{reversibilidade} (adiante, \S\ref{sec:peano}). É por isso que o Monte Carlo não é aqui
uma aproximação inferior à contagem: é a \emph{outra coordenada} da medida, a que o furo
(\S\ref{sec:furar}) obriga.

\medido{um único bit aleatório enche o reticulado $41\times41$ --- o rasto cobre as $1681$ casas,
cobertura $100\%$; a medida de Lebesgue do disco pelo rasto (Metropolis) bate com a contagem ordenada,
e o erro médio sobre $24$ sementes desce de $\sim\!3{,}3\%$ ($20$k passos) para $\sim\!0{,}7\%$ ($800$k),
a convergir; e o par: a curva simples do gato cobre $1/N\to0$ (medida zero) e o rasto aleatório cobre
tudo (medida cheia) (\code{tests/bit\_metropolis.c}).}

\subsection*{A curva de Peano e o acaso: duas faces da mesma moeda}\label{sec:peano}

\textbf{E aqui entra a construção de Gentil, e desfaz uma meia-verdade.} A \emph{curva de Peano} é um
traço \emph{determinista} e contínuo --- 1D --- cuja imagem cobre o quadrado inteiro: \textbf{enche}, de
medida cheia. Logo ``determinista $\Rightarrow$ medida zero'' vale para uma curva \emph{simples}, não
para todas. \textbf{As duas enchem}: o acaso (Metropolis) e a curva de Peano. A questão não é \emph{qual}
enche --- é a \textbf{dualidade} entre elas, e ela é o teorema central.

\begin{teorema}[a Peano e o acaso são os dois lados do teorema central]\label{thm:peano}
As duas enchem o espaço, e são \textbf{duais}. A \textbf{curva de Peano} enche pela \emph{ordem} ---
\textbf{contar}: um índice que visita cada casa uma vez (bijecção), por passos unitários (é curva), o
lado discreto (Hurwitz). O \textbf{rasto de Metropolis} enche pelo \emph{acaso} --- \textbf{medir}: a
fracção do tempo em cada conjunto, o lado contínuo (Gentil). E dão a \emph{mesma} medida:
$\text{contar}=\text{integrar}$, resíduo $0$ (\S\ref{sec:central}). É o \textbf{casamento},
$\mathrm{fp}=1$ --- \emph{duas faces da mesma moeda}.
\end{teorema}

\noindent \textbf{É o teorema central lido no espaço-que-se-enche.} A Peano conta as casas em ordem; o
acaso mede-as pela fracção; e o central diz que a \emph{contagem ordenada} e a \emph{medida aleatória}
são a mesma coisa. Encher pela ordem ou pelo acaso é a estaca e a cruz do mesmo bit --- o par
Gentil\,$\leftrightarrow$\,Hurwitz, casado. E a curva que enche o plano tem uma \textbf{base}, e a base
é o número de níveis do inversor: \textbf{Takahashi} (1986) definiu o controlo directo de torque a
\emph{dois} níveis --- base dois ---, e a extensão a \emph{quatro} níveis e a multinível (a
\emph{monografia}) leva-o a base $N$. A curva de Peano, base \emph{três}, é o caso do trial
$\{-1,0,+1\}$ --- o inversor de três níveis (\code{tests/inversor\_trial.c}) ---, e a sequência de
comutação do DTC \emph{é} essa curva a encher o espaço vectorial; o multinível é a mesma curva noutra
base.

\begin{corolario}[o que se enche é um fecho convexo, e a sua dual é o polar]\label{cor:convexo}
O espaço que a curva enche não é um conjunto qualquer: é um \textbf{fecho convexo}. A modulação --- a
\emph{média} da sequência de comutação --- é a \textbf{combinação convexa} dos vectores do inversor,
logo o alcançável é exactamente o fecho convexo $K$ desses vectores (o \emph{hexágono}, a dois níveis;
um polítopo mais fino, a multinível). A sua \textbf{dual} é o \textbf{polar} $K^{*}$, e fecha por
involução: $K^{**}=K$ (o teorema do bipolar de Rockafellar). \emph{Encher} $K$ (a combinação convexa, a
Peano) e \emph{contar} as suas faces por $K^{*}$ (o polar) são os dois lados do teorema central, agora
em convexidade --- e $\nu\circ\nu=\mathrm{id}$ é o ir-e-voltar $K\mapsto K^{*}\mapsto K$
(\code{tests/polar.c}).
\end{corolario}

\noindent \textbf{E é aqui que se sai da análise e se entra na geometria.} O fecho convexo funda-a: a
medida do encher passa a ser o \emph{volume} de $K$, e a contagem das restrições passa a ser o polar
$K^{*}$. O dual analítico --- \emph{contar} $\leftrightarrow$ \emph{integrar} --- e o dual geométrico
--- $K\leftrightarrow K^{*}$, com $K^{**}=K$ --- são \textbf{o mesmo} $\nu\circ\nu=\mathrm{id}$, lido em
dois eixos. A curva de Peano foi a última coisa da análise; o fecho convexo é a primeira da geometria,
e o teorema do bipolar é o teorema central com outro nome (a face geométrica da dualidade, no
\emph{Catálogo}).

\smallskip
\noindent \textbf{E são as duas leis.} O \textbf{teorema central} --- a involução $\nu\circ\nu=
\mathrm{id}$, contar $\leftrightarrow$ integrar --- é a realização da \textbf{Lei 1} (\emph{a unidade é
dual}); o \textbf{fecho convexo} --- o bipolar $K^{**}=K$ --- é a da \textbf{Lei 2} (\emph{a dualidade é
dual}). E é por isto que os problemas do milénio se partem em duas famílias, pela lei em que cada um
assenta (\emph{Catálogo}, \S\ref{sec:clay}): Lei 1 as reflexões (Riemann, a Poincaré-dualidade) e Lei 2
os bipolares (BSD, Navier--Stokes, Hodge) --- as duas na \emph{fronteira} do fecho ---, e no
\textbf{centro}, o ponto fixo $0$, o Yang--Mills e o $P$ vs $NP$ (a origem e o fim, os dois nulos
colapsados).

\medido{a curva de Peano em $81\times81$ (índice ternário, sem um \code{double}): \textbf{enche} ---
$6561$ casas cobertas, cada uma uma vez, bijecção; é \textbf{curva} --- $0$ passos não-unitários; e a
\textbf{dual fecha} --- contar$\circ$encher $=\mathrm{id}$, resíduo $0$ (\code{tests/peano\_dual.c}); e a
base três é o trial do inversor de três níveis (\code{tests/bit\_metropolis.c} para o acaso, o outro
lado).}

\section{A integral pelo dual: cortar a imagem}\label{sec:lebesgue}

Falta a integral, e ela deriva-se de uma involução já provada em vez de se construir de novo.

\smallskip
\noindent Somar $f$ tem dois cortes possíveis, e são um par:

\begin{center}\small
\begin{tabular}{@{}lll@{}}
\toprule
 & o corte & a pergunta \\
\midrule
pelo \textbf{domínio} & fatias verticais & \emph{onde} estou, e quanto vale aqui \\
pela \textbf{imagem} & fatias horizontais & \emph{quanto} vale, e que tamanho tem o onde \\
\bottomrule
\end{tabular}
\end{center}

\noindent O segundo pede, para cada altura $y$, a \emph{medida} do conjunto $\{f>y\}$ --- e é por
isso que ele só existe depois das secções anteriores: \textbf{cortar a imagem é integrar a medida
dos níveis}. E os dois cortes não são duas teorias: são \emph{o mesmo objecto pelos dois lados da
involução} $f\leftrightarrow f^{-1}$ --- trocar domínio com imagem \emph{é} trocar $f$ pela
inversa, a Lei~2 no nível um. A conservação que já está medida nos metais,
\[
  \int_a^{b} f \;+\; \int_{f(a)}^{f(b)} f^{-1} \;=\; b\,f(b)-a\,f(a),
\]
é exactamente esta troca a fechar sem resto: \emph{o rectângulo reparte-se pelos dois cortes e não
sobra nada}. O clássico chama ao corte da imagem \emph{Lebesgue} e ao do domínio \emph{Riemann};
aqui são as duas coordenadas do mesmo par, e a que corta a imagem é a que herda tudo o que a
medida já provou:

\begin{itemize}
\item \textbf{os furos são de graça}: mudar $f$ num conjunto de medida nula não muda nenhum
$\{f>y\}$ em mais que esse nulo --- a integral não vê o furo. A função que vale $1$ nos racionais
e $0$ no resto integra \textbf{zero}: fura-se $\mathbb{Q}$ (\S\ref{sec:furar}) e o nível ficou
vazio. O corte do domínio \emph{não} a soma --- toda fatia vertical contém os dois valores, e o
par superior/inferior não fecha: \emph{é a assinatura de quem cortou do lado errado do dual};
\item \textbf{a aditividade contável é a do contador}: os pesos somam-se marca a marca, e a série
da metade (Teorema~\ref{thm:furos}) é o que fecha as uniões contáveis;
\item \textbf{a invariância é o relógio}: transladar não mexe no número de marcas --- o pente não
distingue as suas marcas entre si, e a medida herda-lhe a indiferença;
\item \textbf{o que não se mede é o que não tem dual}: um conjunto entra na teoria quando o par
interior/exterior fecha nele --- as duas metades outra vez. Onde as duas coordenadas não colam,
não há medida a atribuir; \emph{não pertence, não por proibição: por não haver como voltar} --- a
segunda consequência do enunciado da \emph{Teoria}, palavra por palavra.
\end{itemize}

\medido{a conservação $\int f+\int f^{-1}=b\,f(b)-a\,f(a)$ está verificada nos três primeiros
metais (na \emph{Teoria}, com o resíduo no truncamento da régua); e no ponto fixo a repartição
separa as duas coordenadas exactamente --- $\int_1^{\varphi}(1+1/x)\,dx=1/\varphi+\ln\varphi$,
com a dualidade a cancelar o transcendente e a deixar a borda. A troca de corte não é analogia:
é uma identidade com medidor.}

\section{O teorema central: contar e integrar são duais}\label{sec:central}

A integral que a secção anterior derivou não é um capítulo a mais: é a \textbf{bijeção} que faltava ao
sistema inteiro. Ela liga os dois lados de que este projecto é feito --- o \emph{discreto} (contar,
inteiro, a álgebra) e o \emph{contínuo} (integrar, a régua com comprimento, a medida) --- e o que os
liga tem nome: a \textbf{estrela}. Primeiro o que é uma bijeção destas.

\begin{definicao}[bijeção dual --- a estrela]\label{def:bijdual}
Uma \textbf{bijeção dual} entre dois objectos $A$ e $B$ é a \emph{interface reversível}: uma
correspondência $\nu$ que é o seu próprio inverso, $\nu\circ\nu=\mathrm{id}$ com resíduo $0$; que
\textbf{liga sem fundir} --- preserva os dois lados em vez de os colar num terceiro; e que \textbf{não
apaga um bit}. É a \textbf{estrela} do corpo estelar: o ponto fixo da involução que troca os dois, a
plena (grau seis, onde soma $=$ produto), o único regime com os dois sentidos.
\end{definicao}

\begin{definicao}[dual]\label{def:dualmed}
$A$ e $B$ são \textbf{duais} quando existe uma bijeção dual entre eles. Nenhum é prévio: a bijeção é o
seu próprio inverso, e o que $A$ é para $B$, $B$ é para $A$. \emph{Não há um melhor} --- dá no mesmo
pela cruz e desce pela estaca.
\end{definicao}

\begin{teorema}[contar e integrar são duais: Gentil e Hurwitz]\label{thm:central}
Sejam o lado \textbf{discreto} --- \emph{Hurwitz}: a norma euclidiana $N(x)=\sum x_i^2$ (a esfera, a
contagem), multiplicativa para o produto \emph{bilinear} exactamente nos graus $1,2,4,8$ --- e o lado
\textbf{contínuo} --- \emph{Gentil}: a norma $\|x\|$ e a fusão homogénea (a hipérbole), sem limite de
grau. Então os dois são \textbf{duais}, e a bijeção dual é a \textbf{estrela}, realizada por esta
teoria da medida: a soma reversível do \S\ref{sec:lebesgue},
\[
  \int_a^b f \;+\; \int_{f(a)}^{f(b)} f^{-1} \;=\; b\,f(b)-a\,f(a),
\]
leva \textbf{contar} ($\sum$, Hurwitz) a \textbf{integrar} ($\int$, Gentil) e a traz de volta. A
bijeção é $f:t\mapsto t^2$ (contar $\leftrightarrow$ raiz), e \emph{nunca se avalia uma raiz}: os dois
cortes --- o do domínio (Riemann, Hurwitz) e o da imagem (Lebesgue, Gentil) --- dão a mesma contagem.
\end{teorema}

\begin{proof}
A bijeção dual é a involução $f\leftrightarrow f^{-1}$, a Lei~2 no nível um, já medida acima
(\S\ref{sec:lebesgue}): o rectângulo reparte-se pelos dois cortes \emph{sem resto}, logo a
correspondência é reversível ($\nu\circ\nu=\mathrm{id}$, resíduo $0$) e é a estrela. Contar (a soma no
domínio) e integrar (o integral na imagem) são os dois lados desta bijeção, logo são duais. E o
\textbf{limite no grau oito é do lado discreto/bilinear} --- Hurwitz classifica o bilinear ---, \emph{não
do objecto}: o lado contínuo, alcançado pela soma reversível e \emph{não por uma raiz}, não o tem. É a
frase desta teoria dita na fronteira álgebra/medida: \textbf{se tens o inteiro de um lado, o integral
leva-o ao contínuo do outro, e volta}.
\end{proof}

\noindent \textbf{E a bijeção é um oscilador --- que é o relógio.} Os dois lados não estão parados:
um \emph{sobe} (a torre das dimensões) enquanto o outro \emph{desce} (a projeção das interpretações),
a estaca troca-os, e o par \textbf{oscila}. Em dois passos, um por lei: a Lei~1 reflecte (período $2$)
e a Lei~2 roda (período $4$). \emph{É o relógio}: um relógio é um oscilador lido pela contagem, um
oscilador é um relógio lido pelo integral --- e a estrela é a bijeção entre os dois. É por isso que
$\pi$ e as constantes contínuas saem por \emph{contagem} nas secções anteriores: conta-se o discreto, e
a inversa veste a roupa contínua no fim.

\smallskip
\noindent \textbf{E a bijeção não é isócrona}: os dois relógios marcam a mesma ordem e correm tempos
diferentes. Combinar o par --- o dual --- custa, no lado discreto (o bilinear), a tábua $n\times n=n^2$
produtos; no grau oito são $8^2=\mathbf{64}$ (\emph{dois lances de um lado, sessenta e quatro do
outro}), e é aí que a norma bilinear fecha e a seguir quebra. \textbf{O oito de Hurwitz é o período
desse relógio}, o tecto do tempo $n^2$ --- não uma lei do mundo. O lado contínuo mede pela norma, um
escalar: outro tempo, sem tecto. A bijeção liga-os, mas não os sincroniza.

\medido{o lado discreto (Hurwitz, o bilinear $\sum x_i^2$) fecha em oito e falha em dezasseis --- o
limite é da contagem, não do objecto; a bijeção dual (a estrela) é reversível: os dois cortes de
$\int f$ dão a mesma contagem e $\int f+\int f^{-1}$ fecha (Young no reticulado), \emph{sem uma raiz};
e o oscilador \emph{é} o relógio --- a Lei~2 roda (período $4$, $a^2+b^2$ imóvel), a Lei~1 reflecte
(período $2$), um sobe enquanto o outro desce e trocam (\code{tests/estelar\_completo.c}); a testemunha
directa do contínuo, $\|xy\|=\|x\|\|y\|$ de $\mathbb{R}^2$ a $\mathbb{R}^7$, está em \code{tests/nne.c}.}

\subsection*{$0{,}\overline9 \ne 1$: a órbita não é o limite --- e por isso não há \emph{malloc}}\label{sec:nove}

\noindent O Teorema~\ref{thm:central} tem um corolário que fecha, e é a Lei~1 (o ponto fixo da involução).
Tome-se a órbita $0{,}9;\,0{,}99;\,0{,}999;\dots$, o $k$-ésimo convergente $\frac{10^k-1}{10^k}$. Do lado
\textbf{discreto} (Hurwitz, \emph{contar}) ela \textbf{nunca chega} a $1$: em todo passo finito o resíduo
é $\frac1{10^k}>0$, com numerador $1$ que nunca é $0$. Do lado \textbf{contínuo} (Gentil, \emph{integrar})
o resíduo $\to 0$, e esse limite \emph{é} o ponto fixo, o real $1$. \textbf{Como órbita, $0{,}\overline9\ne 1$;
como limite, $0{,}\overline9 = 1$} --- os dois lados da mesma bijeção dual, a estrela.

\begin{corolario}[a definição de limite, a continuidade, e a cegueira do contínuo]\label{cor:nove}
A \emph{definição} de limite deriva-se, não se postula: dado qualquer $\varepsilon=\frac1{10^m}>0$,
\textbf{toma-se} $N=m$, e para todo $k>N$ vale $\frac1{10^k}<\varepsilon$ (pois $k>m\Rightarrow 10^k>10^m$)
--- o $\varepsilon$--$N$ clássico, com o $N$ \textbf{construído} a partir de $\varepsilon$. E a
\textbf{continuidade} é comutar com esse limite: para $f$ contínua, $f(1)=\lim f(\text{órbita})$. Uma $f$
linear comuta (o vão $\frac{a}{10^k}\to 0$); mas o \textbf{degrau} $g(x)=[x<1]$ \emph{não}: $g(\text{órbita}_k)=1$
sempre, $g(1)=0$. \textbf{O degrau descontínuo é o único que separa $0{,}\overline9$ de $1$} --- e é o lado
discreto (contar). O contínuo é \textbf{cego} à distinção: é \textbf{Riemann--Lebesgue} --- a média (integrar,
Fourier) manda o resíduo por passo a $0$, dá $1$ para a órbita \emph{e} para a constante $1;1;1;\dots$, e por
isso \textbf{Fourier não mede isto}.
\end{corolario}

\noindent \textbf{E é esta a razão de não haver \emph{malloc}.} Um \code{malloc} é a órbita-acumulação:
guardar cada passo para \emph{chegar} ao limite por soma. Mas a órbita \textbf{nunca fecha} (o resíduo é
$>0$ em todo passo), logo o \emph{buffer} nunca pára de crescer. O ponto fixo --- o real $1$ --- atinge-se
\emph{sem} acumular: pela involução, que é \textbf{reversível} ($\nu\circ\nu=\mathrm{id}$, resíduo $0$, na
prova do Teorema~\ref{thm:central}) e por isso \textbf{não guarda nada} --- um \emph{buffer} é a cópia que a
estrela existe para não fazer (a doutrina «sem memória, e sem estado» do \emph{corpo estelar}). Contar a
órbita numa RAM que cresce é o \code{malloc} que não fecha; \emph{ser} o ponto fixo é a estrela, resíduo $0$,
memória nenhuma. A proibição do \code{malloc} não é higiene de código: é o Teorema~\ref{thm:central} lido do
lado certo --- o limite não se \emph{armazena}, \emph{é-se}.

\medido{$0{,}\overline9$ é o dual de $1$: a órbita $0{,}9;0{,}99;0{,}999;\dots$ nunca chega (resíduo
$\frac1{10^k}$, numerador $1$, nunca $0$, em todo passo finito), e o limite é $1$ (resíduo $\to 0$, o ponto
fixo); a definição de limite deriva-se ($\varepsilon=\frac1{10^m}\Rightarrow N=m$, construtivo), a
continuidade é comutar com o limite e só o degrau descontínuo $[x<1]$ separa os dois lados, enquanto a média
(Riemann--Lebesgue) dá $1$ para ambos --- por isso Fourier não mede a distinção, e por isso a órbita-acumulação
(o \code{malloc}) não fecha e o ponto fixo (a involução, sem memória) fecha: \code{tests/nove\_periodico.c}.}

\section{A medida da espécie: o crânio, não as células}\label{sec:especie}

As secções anteriores medem \emph{conjuntos}. Falta o outro objecto, e ele é o que a cruz sempre
disse: a coordenada que \textbf{mede} é aditiva e dá o tamanho; a que \textbf{ordena} é
multiplicativa e dá a \emph{espécie}. Este parágrafo é a segunda metade, e sem ela o documento
estaria pela metade com o nome do par.

\subsection*{Célula a célula não mede}

Quem quer a espécie de um corpo e o percorre célula a célula não obtém uma medida grosseira:
\textbf{não obtém medida nenhuma}. A espécie não mora nas células --- mora no \emph{crânio}, e
crânios são \textbf{meia dúzia}: classes inteiras, finitas por nível, das quais o corpo ocupa
exactamente uma. O erro não é de resolução, é de categoria --- e por isso mais resolução não o
corrige: quem conta células conta \emph{outra} grandeza.

\subsection*{E «espécie» não é palavra nova: é a régua, com mais um nome}

A \emph{Teoria} abre com os dois grupos de nomes, e este cai no segundo: \textbf{assinatura,
cifra, régua, grau, métrica, fase} --- números que distinguem um corpo dos outros, todos a nomear
\emph{o mesmo caminho na matriz do relógio}. «Espécie» é só mais um sinónimo desse lado, usado
aqui uma vez pela imagem do crânio --- e devolvido: \textbf{o termo canónico continua a ser a
régua}, e é ele que o resto desta secção usa.

\subsection*{O crânio tem dois ossos: a régua, e a dinâmica temporal}

Pelo Teorema do espectro, um corpo \emph{é} exactamente duas coisas --- e por isso o crânio tem
duas colunas, e não uma:

\begin{center}\small
\begin{tabular}{@{}lll@{}}
\toprule
o osso & o que conta & quantas classes \\
\midrule
\multicolumn{3}{@{}l}{\emph{a régua (a leitura):}} \\
\quad a assinatura $(p,q,r)$ & eixos por estado do trial --- cada eixo um relógio & finitas por grau \\
\quad o grau $n=p+q+r$ & quantos relógios & um inteiro \\
\quad o sinal de $\Delta$ & a geometria: elíptica, parabólica, hiperbólica & \textbf{três}, e não há quarta \\
\addlinespace
\multicolumn{3}{@{}l}{\emph{a dinâmica temporal (o que corre):}} \\
\quad o gerador & quem move o corpo, $T\dg=-T$ & um por corpo \\
\quad o período & fecha ou não fecha & \textbf{finito} (a órbita volta) ou \emph{sem retorno} \\
\bottomrule
\end{tabular}
\end{center}

\noindent \textbf{Medir a classe é contar as dimensões inteiras que o corpo ocupa --- e o
relógio dá isso em por-unidade --- e ler que dinâmica o corre.} Não se percorre o corpo:
pergunta-se quantos eixos, em que estados, com que sinal, movidos por que gerador e com que
período. Meia dúzia de números, todos inteiros, e a classe está lida --- como o crânio decide a
espécie sem que ninguém conte uma célula. \emph{Onde o período é finito a órbita fecha e o corpo
tem relógio próprio; onde não é, o corpo tem tempo e não tem relógio --- corre e não volta} ---
e essa é a diferença entre um corpo elíptico e um hiperbólico, dita sem métrica.

\subsection*{E por que Fourier não mede este corpo}

Fourier avalia no círculo --- e \emph{as raízes do metal não estão no círculo}: são recíprocas
pela alfândega, $|\sigma||\sigma'|=1$, uma dentro e uma fora --- a hipérbole. Encostar-lhe o
círculo é medir com o crânio de \textbf{outra} espécie: os coeficientes saem todos não nulos,
nenhuma classe estabiliza, e o factor que não se elimina é o preço da régua alheia. É célula a
célula numa base de fora --- \emph{não-medida}, e agora com o mecanismo dito.

\smallskip
\noindent A transformada que mede é a \textbf{universal}: avaliação nas raízes \emph{da própria
borda} --- os caracteres são a órbita da torção do próprio corpo, e \textbf{trocar o grupo troca a
transformada sem trocar a conta}. Ela não se define: \emph{cai do dual}. E é unitária,
\[
  \|x\|^{2} \;=\; \|\mathcal{F}x\|^{2},
\]
que é a frase desta teoria no espectro: \textbf{nenhuma medida vaza na decomposição} --- os
pedaços somam o objecto sem sobra e sem excesso, que é o Princípio da Medida Consistente lido nos
coeficientes.

\subsection*{E a roupa da espiral veste-se no fim, pela inversa}

A ordem é a do corpo estelar, e não se negoceia: \textbf{o esqueleto inteiro primeiro, a roupa no
fim}. Conta-se a espécie em por-unidade --- os inteiros do crânio --- e só então a
\emph{transformada inversa universal} reconstrói o corpo contínuo do seu espectro:
$\mathcal{F}^{-1}\mathcal{F}=\mathrm{id}$, a volta exacta, com o período quatro do esquilo
($\mathcal{F}^{2}$ é a reflexão, $\mathcal{F}^{4}$ a identidade). \emph{A espiral é a roupa} ---
a forma contínua que se veste sobre a contagem --- e quem começa pela roupa mede a roupa: dígitos
da forma em vez dos inteiros do crânio, células em vez de espécie. \textbf{A medida é a contagem;
a espiral é o que a inversa devolve depois de contado.}

\medido{a ida-e-volta da transformada universal é exacta nos $121$ elementos do anel, resíduo
$0$ e \emph{sem normalização nenhuma} --- $V^{-1}(V(a))=a$ ---, e ela diagonaliza:
$\mathcal{F}(a\circledast b)=\mathcal{F}(a)\mathcal{F}(b)$, resíduo $0$
(\code{tests/universal.c}); a definição com $1/\sqrt n$ em cada lado fecha no ponto
$\sum_k\chi_k(j)\chi_{-k}(j')=n\,\delta_{j,j'}$ (\code{tests/transformada.c});
$\mathcal{F}^{-1}\mathcal{F}(D)=D$ para qualquer $D$ na máquina completa
(\code{tests/completo.c}); e a assinatura conta: três eixos dão $8$ estados, as seis permutações
preservam o peso, e o controlo com dois eixos não dá trialidade
(\code{tests/trialidade.c}).}

\section{A velocidade: por que contar é rápido}

A queixa que originou este documento é operacional, e a resposta também é.

\begin{center}\small
\begin{tabular}{@{}llll@{}}
\toprule
o método & o gesto & o custo & o que devolve \\
\midrule
aproximar & dígito a dígito & um passo por dígito, \emph{para sempre} & meia coordenada \\
tentar & propor e filtrar & o acaso decide o ritmo & o mesmo número, pior \\
\textbf{contar} & marcas inteiras numa escala & \emph{uma} varredura & o par, enquadrado \\
\textbf{furar} & cobrir o contável & um encaixante por ponto & o resto, inteiro \\
\bottomrule
\end{tabular}
\end{center}

\noindent O número está medido no \code{pimonte}: a contagem com $R=3000$ chega a $15$
milionésimos onde o método térmico fica em $84$ --- \emph{e o aproximador decimal nem entra na
tabela do medidor, porque nunca fecha}: cada dígito novo é um passo novo, e não há passo em que a
metade que falta chegue. \textbf{Contar fecha por escala; furar fecha por ponto; aproximar não
fecha.} É a mesma economia do medidor-zero: resolver e provar na mesma passagem, em vez de
resolver para sempre e provar nunca.

\section{A medida da própria medição}

O fecho é o do corpo estelar, agora com o vocabulário completo. Uma asserção que compara contra um
valor escrito é um \emph{aproximador}: encosta por um lado, e passa no objecto certo e no trocado.
O medidor-zero é um \emph{encaixante}: a reversão dá as duas coordenadas na mesma passagem, e o
resíduo é o tamanho do que ficou entre elas --- \textbf{medir é encaixar, e o zero é o encaixe que
fechou}. E a bateria inteira é a medida de contagem sobre os medidores: verdes contados um a um,
com o total a acusar quem desaparece --- \emph{o furo que ela não perdoa é exactamente o que esta
teoria diz que não pesa nada e por isso ninguém vê sair}.

\medido{a promoção $S=(x+x\dg)/2$, $A=(x-x\dg)/2$ reconstrói $x$ em $48\,861$ casos com nenhuma
das $97\,722$ somas ímpar --- a metade é exacta porque a involução a garante
(\code{tests/promove.c}); a prova dos nove fecha em $23\,200$ pares e acusa o erro de uma unidade
em todos (\code{tests/reversao.c}); e pela porta única da ISA a máquina fabrica as suas
constantes --- o $3$ do \code{CMP} dos zeros, o $2$ dos iguais, o $1$ do \code{SUB} --- e mede-se
a si própria pelo reflexo, com o resíduo a sair pela própria porta (\code{tests/isa\_wasm.js}).}

\section*{A dívida, paga}

A primeira versão deste documento nomeava dois medidores por existir. Existem, e correm na
bateria como \code{tests/medida.c}, tudo em inteiros e por produto cruzado.

\medido{$\lambda([a,b])=b-a$ \emph{exacto} em $588$ intervalos racionais, marca a marca contra a
forma fechada, em quatro escalas comensuráveis cada --- e o controlo: fora da escala comensurável
o contador \emph{enquadra} a menos de um passo em $75$ de $75$ e só iguala em $10$: a exactidão é
da escala, não do acaso. O furo: na escala $q=20\,011$, cobrir os racionais enumerados com
encaixantes de $\varepsilon/2^{k}$ ($\varepsilon=1/10$) deixa $18\,760$ marcas livres ---
\emph{a régua não se move} --- e o controlo é a série errada, $\varepsilon/k$, que passa de
$\varepsilon$ e passa do intervalo inteiro. Dirichlet: o corte do domínio \textbf{não fecha em
escala nenhuma} (o gap superior$-$inferior é $1$ em todas, com o testemunho irracional a sair da
\emph{borda} --- $0<\sigma-1<1$, e $\pm1$ negados pela raiz racional --- sem avaliar um decimal);
o corte da imagem fecha em \textbf{zero}, com a soma termo a termo a bater com a forma fechada
$\varepsilon(2^{n}-1)/2^{n}$ e a ficar abaixo de cada $\varepsilon$. Os dois cortes: colunas
$=$ linhas em $177$ pares (grelha, função), transposição exacta, e o rectângulo reparte-se sem
resto --- $\int f+\int f^{-1}=$ área, na grelha inteira. E a mutação: furar \emph{um} ponto custa
\emph{uma} marca, em todas as escalas (\code{tests/medida.c}).}

\vfill
{\footnotesize\color{cinza}\noindent
Teoria, catálogo, medidores e o resto do sistema em
\href{https://goldenkingdom.patriatechnology.com}{goldenkingdom.patriatechnology.com}.
CC BY 4.0 (textos) $\cdot$ MIT (código).\par}

\end{document}
